
\subsection{勾股恒等式}
\label{subsec:pythagorean}



\begin{figure}[!h]
    \centering
    \subfloat[勾股定理 \href{https://socratic.org/questions/58a20042b72cff17846a6feb}{（来源）}]{
  \includegraphics[width=0.3\linewidth]{figures/pythagorean-theorem.png}

  \label{fig:pythagorean-theorem}
}
    \subfloat[勾股恒等式 \href{https://socratic.org/questions/58a20042b72cff17846a6feb}{（来源）}]{
  \includegraphics[width=0.3\linewidth]{figures/pythagorean-identity.png}

  \label{fig:pythagorean-identity}
}
  \caption{}
\end{figure}

\begin{tcolorbox}[title={\textbf{\tboxtheorem{\ref*{subsec:pythagorean}} 勾股恒等式}}]
$\sin^2 \theta + \cos^2 \theta = 1$
\end{tcolorbox}

\begin{proof}
$ $
\begin{enumerate}

\item 在 \autoref{fig:pythagorean-theorem} 中，根据勾股定理，有 $a^2 + b^2 = c^2$。

\item 在 \autoref{fig:pythagorean-identity} 中，根据正弦与余弦的定义，

$\sin\theta = \dfrac{a}{c}, \text{ } \cos\theta = \dfrac{b}{c}$

\item 结合步骤 1 和步骤 2：

$\sin^2 \theta + \cos^2 \theta = \dfrac{a^2}{c^2} + \dfrac{b^2}{c^2} = \dfrac{a^2 + b^2}{c^2} = 1$


\end{enumerate}
\end{proof}

\subsection{虚数}
\label{subsec:imaginary}


\begin{tcolorbox}[title={\textbf{\tboxdef{\ref*{subsec:imaginary}} 虚数}}]

\begin{itemize}

\item 将 \boldmath{$i$} 定义为满足 $i^2 =-1$ 的虚数单位。
\item \textbf{复数}是形如 $a + b\,i$ 的数，其中 $a$ 和 $b$ 均为实数（例如，$a = 5.6 + 4.3i)$

\item \textbf{实数}不包含虚部（例如，$a = 13.4$）。

\item 若 $a$ 和 $\overline{a}$ 具有相同的实部以及符号相反的虚部，则 $\overline{a}$ 是 $a$ 的\textbf{共轭}。

(e.g., $a = 3 + i\cdot3.4$, \text{ } $\overline{a} = 3 - i\cdot3.4$)

\item \textbf{Hermitian 向量}指后半部分元素为前半部分元素逆序排列后复共轭的向量，如下面的 $n$ 维向量所示：

$\vec{v} = (v_1, v_2, v_3, \cdots, v_{\frac{n}{2}-1}, v_{\frac{n}{2}}, \overline{v}_{\frac{n}{2}}, \overline{v}_{\frac{n}{2} - 1}, \cdots, \overline{v}_3, \overline{v}_2, \overline{v}_1)$


\end{itemize}
\end{tcolorbox}



\subsection{欧拉公式}
\label{subsec:euler}




\begin{tcolorbox}[title={\textbf{\tboxdef{\ref*{subsec:euler}} 欧拉公式}}]

$e^{i\theta} = \cos\theta + i\cdot\sin\theta$

\end{tcolorbox}

\begin{figure}[!h]
    \centering
  \includegraphics[width=0.4\linewidth]{figures/euler-formula.png}
  \caption{复平面单位圆上的欧拉公式示意图。\href{https://en.wikipedia.org/wiki/Euler's_formula}{（来源）}}
\end{figure}

在 \autoref{fig:complex-plane} 中，$e^{i\theta}$ 的值表示为复平面圆上的一个坐标点，其中 $x$ 轴表示实部，$y$ 轴表示虚部。随着 $\theta$ 增大，虚部（$\sin\theta$）在 $1$ 与 $-1$ 之间振荡，在虚轴上分别达到 $i$ 和 $-i$；实部（$\cos\theta$）在 $1$ 与 $-1$ 之间振荡，周期为 $2\pi$。


\subsection{以复数上分圆多项式的根构造 Vandermonde 矩阵}
\label{subsec:vandermonde-euler}

 
 
本小节利用复数上第 $\mu$ 个分圆多项式的 $n$ 个互异根，构造 Vandermonde 矩阵（见 \autoref{subsec:vandermonde}），其中 $\mu$ 是 2 的幂，具体如下：



\begin{tcolorbox}[title={\textbf{\tboxtheorem{\ref*{subsec:vandermonde-euler}} 以复数上指标为 2 的幂的分圆多项式的根构造 Vandermonde 矩阵}}]

假设有一个 $n \times n$ 的 Vandermonde 矩阵（其中 $n$ 是 2 的幂），由第 $\mu$ 个分圆多项式的 $n$ 个互异根组成（见 \autoref{subsec:cyclotomic-theorem} 中的定理~\ref*{subsec:cyclotomic-theorem}.1），其中 $\mu$ 是 2 的幂，且 $n = \dfrac{\mu}{2}$。换言之，$V = \mathit{Vander}(x_0, x_1, \cdots, x_{n-1})$，并且对于 $1 \leq j \leq n$，每个根均为 $x_j = (e^{i\pi/n})^{2j-1}$，即 $\mu$ 次本原单位根。则有：

$V \cdot V^T = \begin{bmatrix}
0 & \cdots & 0 & 0 & n\\
0 & \cdots & 0 & n & 0\\
0 & \cdots & n & 0 & 0\\
\vdots & \iddots & \vdots & \vdots & \vdots \\
n & 0 & 0 & \cdots & 0\\
\end{bmatrix} = n \cdot I^R_n$

$ $

并且 $V^{-1} = \dfrac{V^T \cdot I_n^R}{n}$


\end{tcolorbox}

\begin{proof}

$ $

\begin{enumerate}

\item 给定 $\omega = e^{i\pi/n}$，其中每个 $x_j = (\omega)^{2j-1}$。因此，可以展开如下：

$V \cdot V^T = \begin{bmatrix}
1 & (\omega) & (\omega)^2 & \cdots & (\omega)^{n-1}\\
1 & (\omega^3) & (\omega^3)^2 & \cdots & (\omega^3)^{n-1}\\
1 & (\omega^5) & (\omega^5)^2 & \cdots & (\omega^5)^{n-1}\\
\vdots & \vdots & \vdots & \ddots & \vdots \\
1 & (\omega^{2n-1}) & (\omega^{2n-1})^2 & \cdots & (\omega^{2n-1})^{n-1}\\
\end{bmatrix} 
\cdot 
\begin{bmatrix}
1 & 1 & 1 & \cdots & 1\\
(\omega) & (\omega^3) & (\omega^5) & \cdots & (\omega^{2n-1})\\
(\omega)^2 & (\omega^3)^2 & (\omega^5)^2 & \cdots & (\omega^{2n-1})^2\\
\vdots & \vdots & \vdots & \ddots & \vdots \\
(\omega)^{n-1} & (\omega^3)^{n-1} & (\omega^5)^{n-1} & \cdots & (\omega^{2n-1})^{n-1}\\
\end{bmatrix} $

$ $


$=
\begin{bmatrix}
\sum\limits_{k=0}^{n-1} \omega^{2k} & \sum\limits_{k=0}^{n-1} \omega^{4k}  & \sum\limits_{k=0}^{n-1} \omega^{6k} & \cdots & \sum\limits_{k=0}^{n-1} \omega^{2nk} \\

\sum\limits_{k=0}^{n-1} \omega^{4k} & \sum\limits_{k=0}^{n-1} \omega^{6k}  & \sum\limits_{k=0}^{n-1} \omega^{8k} & \cdots & \sum\limits_{k=0}^{n-1} \omega^{2k(n+1)} \\

\sum\limits_{k=0}^{n-1} \omega^{6k} & \sum\limits_{k=0}^{n-1} \omega^{8k} & \sum\limits_{k=0}^{n-1} \omega^{10k} & \cdots & \sum\limits_{k=0}^{n-1} \omega^{2k(n+2)} \\

\vdots & \vdots & \vdots & \ddots & \vdots \\
\sum\limits_{k=0}^{n-1} \omega^{2nk} & \sum\limits_{k=0}^{n-1} \omega^{2(n+1)k} & \sum\limits_{k=0}^{n-1} \omega^{2(n+2)k} & \cdots & \sum\limits_{k=0}^{n-1} \omega^{2(n+n-1)k} \\

\end{bmatrix}$

$ $

\item 矩阵 $V \cdot V^T$ 的副对角线元素为 $\sum\limits_{k=0}^{n-1} \omega^{2nk}$。可以推导：

$\sum\limits_{k=0}^{n-1} \omega^{2nk} = \sum\limits_{k=0}^{n-1} (e^{i\pi/n})^{2nk} = \sum\limits_{k=0}^{n-1}e^{2\pi k i} = \sum\limits_{k=0}^{n-1}(\cos(2\pi k) + i\sin(2\pi k)) = \sum\limits_{k=0}^{n-1} (1 + 0) = n$

这说明，矩阵 $V \cdot V^T$ 的副对角线元素为 $n$。

$ $

\item 接下来，证明矩阵 $V \cdot V^T$ 除副对角线外的所有元素均为 0。也就是说，证明：

$\sum\limits_{k=0}^{n-1} \omega^{2k} = \sum\limits_{k=0}^{n-1} \omega^{4k} = \sum\limits_{k=0}^{n-1} \omega^{6k} = \cdots = \sum\limits_{k=0}^{n-1} \omega^{2(n-1)k} = \sum\limits_{k=0}^{n-1} \omega^{2(n+1)k} = \cdots = \sum\limits_{k=0}^{n-1} \omega^{2(2n-1)k} = 0$

证明中需要使用等比数列求和公式 $\sum\limits_{i=0}^{n-1}x^i = \dfrac{x^n - 1}{x - 1}$：

\begin{tcolorbox}[title={\textbf{\tboxtheorem{\ref*{subsec:vandermonde-euler}.1} 等比数列求和公式}}]

设等比数列之和为 $S_n = 1 + x + x^2 + \cdots + x^{n - 1}$

则 $x \cdot S_n = x + x^2 + x^3 + \cdots + x^{n}$

$x \cdot S_n - S_n = (x + x^2 + x^3 + \cdots + x^{n}) - (1 + x + x^2 + \cdots + x^{n - 1}) = x^n - 1$

$S_n\cdot (x - 1) = x^n - 1$

$S_n = \dfrac{x^n - 1}{x - 1}$ \textcolor{red}{ $\rhd$ 其中要求 $x \neq 1$}

\end{tcolorbox}

利用等比数列求和公式 $\sum\limits_{i=0}^{n-1}x^i = \dfrac{x^n - 1}{x - 1}$，

对于 $1 \leq m \leq (2n - 1)$，有 $\sum\limits_{k=0}^{n-1} \omega^{2mk} = \dfrac{(\omega^{2m})^n - 1}{\omega^{2m} - 1} = \dfrac{(\omega^{2n})^{m} - 1}{\omega^{2m} - 1} = \dfrac{1 - 1}{\omega^{2m} - 1} = 0$ \textcolor{red}{ $\rhd$ 因为 $\textsf{Ord}(\omega) = 2n$}

$ $

因此，

$\sum\limits_{k=0}^{n-1} \omega^{2k} = \sum\limits_{k=0}^{n-1} \omega^{4k} = \sum\limits_{k=0}^{n-1} \omega^{6k} = \cdots = \sum\limits_{k=0}^{n-1} \omega^{2(n-1)k} = \sum\limits_{k=0}^{n-1} \omega^{2(n+1)k} = \cdots = \sum\limits_{k=0}^{n-1} \omega^{2(2n-1)k} = 0$

$ $

\item 根据步骤 2 和步骤 3 的证明，$V \cdot V^T = 
\begin{bmatrix}
0 & \cdots & 0 & 0 & n\\
0 & \cdots & 0 & n & 0\\
0 & \cdots & n & 0 & 0\\
\vdots & \iddots & \vdots & \vdots & \vdots \\
n & 0 & 0 & \cdots & 0\\
\end{bmatrix} = n \cdot I_n^R$

\item 给定 $V \cdot V^T = n \cdot I_n^R$，

$V^{-1} \cdot V \cdot V^T = V^{-1} \cdot n \cdot I_n^R$

$V^T = V^{-1} \cdot n \cdot I_n^R$

$V^T \cdot I_n^R = V^{-1} \cdot n \cdot I_n^R  \cdot I_n^R$

$V^T \cdot I_n^R = V^{-1} \cdot n$ \textcolor{red}{\text{ } \# 因为 $I_n^R  \cdot I_n^R = I_n$}

$V^{-1} = \dfrac{V^T \cdot I_n^R}{n}$

\end{enumerate}
\end{proof}

后续在 CKKS 方案（\autoref{sec:ckks}）中，将使用 $V^{-1}$ 将复向量编码为实向量，并使用 $V^T$ 将实向量解码为复向量（见 \autoref{subsec:ckks-enc-dec}）。

$ $

%\para{Condition for $\bm \mu$:} It's worthwhile to note that the property $V\cdot V^T = n\cdot I_n^R$ does not hold if $\mu$ (denoting the $\mu$-th cyclotomic polynomial) is not a power of 2. In particular, step 3 of the proof does not hold anymore if $\mu$ is not a power of 2:

%$\sum\limits_{k=0}^{n-1} \omega^{2k} \neq \sum\limits_{k=0}^{n-1} \omega^{4k} \neq \sum\limits_{k=0}^{n-1} \omega^{6k} \neq \cdots \neq \sum\limits_{k=0}^{n-1} \omega^{2(n-1)k} \neq \sum\limits_{k=0}^{n-1} \omega^{2(n+1)k} \neq \cdots \neq \sum\limits_{k=0}^{n-1} \omega^{2(2n-1)k} \neq 0$


\para{关于 $\bm \mu$ 的条件：}恒等式 $V \cdot V^T = n \cdot I_n^R$ 只依赖于
$\textsf{Ord}(\omega) = 2n$（步骤 3），因此对任意 $\mu = 2n$（即任意偶数）都成立。
同样，节点 $\omega, \omega^3, \ldots, \omega^{2n-1}$ 是 $X^n + 1$ 的 $n$ 个根，
这对任意 $n$ 都成立，因为 $(\omega^{2j-1})^n = \omega^{2nj}\cdot\omega^{-n} = 1\cdot(-1) = -1$。

不过，将 $\mu$ 取为 2 的幂的原因并非来自本定理，
而是来自全同态加密方案本身。在实际应用中，CKKS、BFV 和 BGV 通常采用
分圆环 $\mathbb{Z}[X]/(X^n + 1)$，其中 $n$ 是 2 的幂。这种选择使 $X^n + 1 = \Phi_{2n}(X)$ 在有理数域上不可约，而这一性质成立
当且仅当 $n$ 是 2 的幂；同时，它还支持高效的负循环 NTT 乘法，
并具有清晰的 Ring-LWE 安全性归约。一旦将
$n$（从而也将 $\mu = 2n$）固定为 2 的幂，这些节点就自动成为本原的
$\mu$ 次单位根：根据定理~\ref*{subsec:roots-theorem}.4（\autoref{subsec:roots-theorem}），
$\omega^k$ 是 $\mu$ 次本原单位根，当且仅当 $\gcd(\mu, k) = 1$；而每个奇数指数
$1, 3, \ldots, 2n-1$ 都与 $\mu = 2n$ 互素，当且仅当 $2n$ 不含奇素因子，即 $2n$
是 2 的幂。因此，$\{\omega, \omega^3, \ldots, \omega^{2n-1}\}$ 恰好为本原的
$(\mu=2n)$ 次单位根，是选择二幂分圆环的结果。同时，只要 $\mu$ 为偶数，$V \cdot V^T = n \cdot I_n^R$ 仍然成立。



\subsection{以环（$\mathbb{Z}_p$）上分圆多项式的根构造 Vandermonde 矩阵}
\label{subsec:vandermonde-euler-integer-ring}

定理~\ref*{subsec:vandermonde-euler}（见 \autoref{subsec:vandermonde-euler}）证明了 $V \cdot V^T = n \cdot I^R_n$，其中 $V$ 为 Vandermonde 矩阵 $V = \mathit{Vander}(x_0, x_1, \cdots, x_{n-1})$，每个 $x_i$ 是 $X \in \mathbb{C}$（即复数）上的 $\mu$ 次本原单位根，且 $\mu$ 为 2 的幂。本小节将证明，即使每个 $x_i$ 是 $X \in \mathbb{Z}_p$（即环）上的 $\mu$ 次本原单位根，关系 $V \cdot V^T = n \cdot I^R_n$ 也成立。具体地，将证明定理~\ref*{subsec:vandermonde-euler}：


\begin{tcolorbox}[title={\textbf{\tboxtheorem{\ref*{subsec:vandermonde-euler-integer-ring}} 以环（$\mathbb{Z}_p$）上指标为 2 的幂的分圆多项式的根构造 Vandermonde 矩阵}}]

证明方法与定理~\ref*{subsec:vandermonde-euler}（见 \autoref{subsec:vandermonde-euler}）相同。假设有一个 $n \times n$ 的 Vandermonde 矩阵（其中 $n$ 是 2 的幂），由环 $X \in \mathbb{Z}_p$ 上第 $\mu$ 个分圆多项式的 $n$ 个互异根组成，其中 $\mu$ 是 2 的幂且 $n = \dfrac{\mu}{2}$。换言之，$V = \mathit{Vander}(x_0, x_1, \cdots, x_{n-1})$，其中每个 $x_i$ 都是 $X^n + 1$ 的根，即 $(\mu=2n)$ 次本原单位根。则有：

$V \cdot V^T = \begin{bmatrix}
0 & \cdots & 0 & 0 & n\\
0 & \cdots & 0 & n & 0\\
0 & \cdots & n & 0 & 0\\
\vdots & \iddots & \vdots & \vdots & \vdots \\
n & 0 & 0 & \cdots & 0\\
\end{bmatrix} = n \cdot I^R_n$

$ $

并且 $V^{-1} = n^{-1}\cdot V^T \cdot I_n^R$


\end{tcolorbox}
\begin{proof}
$ $
\begin{enumerate}
\item 将 $V \cdot V^T$ 展开如下：

$V \cdot V^T = \begin{bmatrix}
1 & (\omega) & (\omega)^2 & \cdots & (\omega)^{n-1}\\
1 & (\omega^3) & (\omega^3)^2 & \cdots & (\omega^3)^{n-1}\\
1 & (\omega^5) & (\omega^5)^2 & \cdots & (\omega^5)^{n-1}\\
\vdots & \vdots & \vdots & \ddots & \vdots \\
1 & (\omega^{2n-1}) & (\omega^{2n-1})^2 & \cdots & (\omega^{2n-1})^{n-1}\\
\end{bmatrix} 
\cdot 
\begin{bmatrix}
1 & 1 & 1 & \cdots & 1\\
(\omega) & (\omega^3) & (\omega^5) & \cdots & (\omega^{2n-1})\\
(\omega)^2 & (\omega^3)^2 & (\omega^5)^2 & \cdots & (\omega^{2n-1})^2\\
\vdots & \vdots & \vdots & \ddots & \vdots \\
(\omega)^{n-1} & (\omega^3)^{n-1} & (\omega^5)^{n-1} & \cdots & (\omega^{2n-1})^{n-1}\\
\end{bmatrix} $

$ $


$=
\begin{bmatrix}
\sum\limits_{k=0}^{n-1} \omega^{2k} & \sum\limits_{k=0}^{n-1} \omega^{4k}  & \sum\limits_{k=0}^{n-1} \omega^{6k} & \cdots & \sum\limits_{k=0}^{n-1} \omega^{2nk} \\

\sum\limits_{k=0}^{n-1} \omega^{4k} & \sum\limits_{k=0}^{n-1} \omega^{6k}  & \sum\limits_{k=0}^{n-1} \omega^{8k} & \cdots & \sum\limits_{k=0}^{n-1} \omega^{2k(n+1)} \\

\sum\limits_{k=0}^{n-1} \omega^{6k} & \sum\limits_{k=0}^{n-1} \omega^{8k} & \sum\limits_{k=0}^{n-1} \omega^{10k} & \cdots & \sum\limits_{k=0}^{n-1} \omega^{2k(n+2)} \\

\vdots & \vdots & \vdots & \ddots & \vdots \\
\sum\limits_{k=0}^{n-1} \omega^{2nk} & \sum\limits_{k=0}^{n-1} \omega^{2(n+1)k} & \sum\limits_{k=0}^{n-1} \omega^{2(n+2)k} & \cdots & \sum\limits_{k=0}^{n-1} \omega^{2(n+n-1)k} \\

\end{bmatrix}$

$ $

其中，$\omega$（即 $(\mu=2n)$ 次本原单位根）的阶为 $2n$。

$ $

\item 矩阵 $V \cdot V^T$ 的副对角线元素为 $\sum\limits_{k=0}^{n-1} \omega^{2nk}$。由于 $\textsf{Ord}_p(\omega) = 2n$，可知 $\omega^{2n} \equiv 1 \bmod p$。因此，矩阵 $V \cdot V^T$ 的每个副对角线元素都为 $\sum\limits_{k=0}^{n-1} 1 = n$。

$ $

\item 接下来，证明矩阵 $V \cdot V^T$ 除副对角线外的其他元素均为 $0$。也就是说，证明：

$\sum\limits_{k=0}^{n-1} \omega^{2k} = \sum\limits_{k=0}^{n-1} \omega^{4k} = \sum\limits_{k=0}^{n-1} \omega^{6k} = \cdots = \sum\limits_{k=0}^{n-1} \omega^{2(n-1)k} = \sum\limits_{k=0}^{n-1} \omega^{2(n+1)k} = \cdots = \sum\limits_{k=0}^{n-1} \omega^{2(2n-1)k} = 0$

$ $

由等比数列求和公式（定理~\ref*{subsec:vandermonde-euler}.1），上式成立。如定理~\ref*{subsec:vandermonde-euler} 的证明步骤 3 所示，每个元素都可写成 $\sum_{k=0}^{n-1} (\omega^{2m})^k$，其中 $m$ 为某个整数，且 $2m$ 不是 $2n$ 的倍数。令 $r = \omega^{2m}$，则有：

$\sum\limits_{k=0}^{n-1} r^k = \dfrac{r^n - 1}{r - 1} = \dfrac{(\omega^{2m})^n - 1}{\omega^{2m} - 1} = \dfrac{(\omega^n)^{2m} - 1}{\omega^{2m} - 1}$

$ $

由于 $\omega$ 是 $X^n + 1$ 的根，所以 $\omega^n \equiv -1 \bmod p$。因此：

$\dfrac{(-1)^{2m} - 1}{\omega^{2m} - 1} = \dfrac{1 - 1}{\omega^{2m} - 1} = 0$

$ $

由此，所有非副对角线位置上的和都为 0。

$ $

\item 根据步骤 1 和步骤 2，矩阵 $V \cdot V^T$ 的副对角线元素为 $n$，其他位置的元素都为 $0$。

\item 现在推导 $V^{-1}$ 的公式。给定 $V \cdot V^T = n \cdot I_n^R$，

$V^{-1} \cdot V \cdot V^T = V^{-1} \cdot n \cdot I_n^R$

$V^T = V^{-1} \cdot n \cdot I_n^R$

$V^T \cdot I_n^R = V^{-1} \cdot n \cdot I_n^R  \cdot I_n^R$

$V^T \cdot I_n^R = V^{-1} \cdot n$ \textcolor{red}{\text{ } \# 因为 $I_n^R  \cdot I_n^R = I_n$}

$ $

这里有一点需要注意：模运算不支持直接进行通常的数值除法（如 \autoref{subsec:modulo-division} 所述）。因此，定理~\ref*{subsec:vandermonde-euler}（见 \autoref{subsec:vandermonde-euler}）中的公式 $V^{-1} = \dfrac{V^T \cdot I_n^R}{n}$ 不能直接用于当前情形，因为模 $p$ 的运算不允许直接用 $n$ 除 $V^T \cdot I_n^R$。所以，改为将 $V^T \cdot I_n^R$ 乘以 $n$ 的逆元，即 $n^{-1}$。继续推导如下：

$V^T \cdot I_n^R = V^{-1} \cdot n$

$V^T \cdot I_n^R \cdot n^{-1}= V^{-1} \cdot n \cdot n^{-1}$

$V^{-1} = n^{-1}\cdot V^T \cdot I_n^R$

\end{enumerate}
\end{proof}

最终证明了 $V\cdot V^T = n \cdot I_n^R$ 和 $V^{-1} = n^{-1}\cdot V^T \cdot I_n^R$。后续在 BFV 方案（\autoref{sec:bfv}）中，将使用 $V^{-1}$ 将整数向量编码为多项式系数向量，再使用 $V^T$ 将其解码回整数向量（见 \autoref{subsec:bfv-batch-encoding}）。

$ $

\para{关于 $\bm \mu$ 的条件：}与 CKKS 中的情况类似，如果 $\mu$（表示第 $\mu$ 个分圆多项式的指标）不是 2 的幂，性质 $V\cdot V^T = n\cdot I_n^R$ 就不成立。具体来说，当 $\mu$ 不是 2 的幂时，证明的步骤 3 不再成立：

$\sum\limits_{k=0}^{n-1} \omega^{2k} \neq \sum\limits_{k=0}^{n-1} \omega^{4k} \neq \sum\limits_{k=0}^{n-1} \omega^{6k} \neq \cdots \neq \sum\limits_{k=0}^{n-1} \omega^{2(n-1)k} \neq \sum\limits_{k=0}^{n-1} \omega^{2(n+1)k} \neq \cdots \neq \sum\limits_{k=0}^{n-1} \omega^{2(2n-1)k} \neq 0$